\documentclass[11pt,a4paper]{article}

\usepackage[utf8]{inputenc}
\usepackage[margin=1in]{geometry}
\usepackage{amsmath,amssymb,amsfonts,mathrsfs}
\usepackage{graphicx}
\usepackage{booktabs}
\usepackage{multirow}
\usepackage{array}
\usepackage{xcolor}
\usepackage{hyperref}
\usepackage{cite}
\usepackage{float}
\usepackage{caption}
\usepackage{subcaption}
\usepackage{url}
\usepackage{microtype}
\usepackage{enumitem}

\hypersetup{
    colorlinks=true,
    linkcolor=blue!70!black,
    citecolor=blue!70!black,
    urlcolor=blue!70!black
}

\title{\textbf{GraphEpiRec: A Heterogeneous Knowledge Graph Neural Network for Clinical Note-Aware Antiseizure Medication Recommendation in Epilepsy}}

\author{
  \textbf{M J Yogesh}$^1$ \quad \textbf{Dr. Karthikeyan J}$^{1,*}$ \\[0.5em]
  $^1$School of Computer Science Engineering and Information Systems,\\
  Vellore Institute of Technology, Vellore 632 014, Tamil Nadu, India \\[0.3em]
  \small{$^*$Corresponding Authors: \href{mailto:yogeshmj.nie@gmail.com}{yogeshmj.nie@gmail.com}, \href{mailto:karthikeyan.jk@vit.ac.in}{karthikeyan.jk@vit.ac.in}}
}

\date{\today}

\begin{document}

\maketitle

\begin{abstract}
Selecting an appropriate antiseizure medication (ASM) regime for an individual with epilepsy requires reasoning over a complex web of clinical factors---seizure type, comorbidities, drug-drug interactions, contraindications, and individual patient tolerance. Traditional clinical decision support and collaborative filtering approaches fail to capture the multi-relational graph structure of electronic health record (EHR) data and clinical text notes. We present \textbf{GraphEpiRec}, a multimodal framework that unifies a Relational Graph Convolutional Network (RGCN) over a heterogeneous EHR knowledge graph with a BioBERT-based clinical note context encoder and a gated multimodal fusion unit. The knowledge graph was constructed from a curated cohort of 3,641 epilepsy patients in MIMIC-IV v3.1 (ICD-9 345.x / ICD-10 G40.x), incorporating 3,787 nodes (patients, ASMs, EpSO-aligned diagnoses, and Elixhauser comorbidities) and 98,347 typed edges across six semantic relation types. To prevent transductive graph data leakage, patient-patient similarity edges ($68,526$) were constructed strictly within training set partitions. Evaluated under a 5-fold stratified cross-validation protocol on an independent test set ($n=547$), GraphEpiRec achieved standard monotonic precision decay ($P@1 = 0.812 \pm 0.014$, $P@3 = 0.645 \pm 0.011$, $P@5 = 0.498 \pm 0.009$), robust recall ($R@5 = 0.852 \pm 0.008$), Hit Ratio ($HR@5 = 0.931 \pm 0.007$), NDCG@10 of $0.791 \pm 0.012$, and Mean Reciprocal Rank (MRR) of $0.834 \pm 0.011$, outperforming state-of-the-art baselines ($p < 0.01$, Wilcoxon signed-rank test). Crucially, explicit contraindication edge propagation and drug-interaction modeling reduced top-5 recommendation safety violations from $4.8\% \pm 0.5\%$ (collaborative filtering) to $1.7\% \pm 0.3\%$. Subgroup evaluations confirmed 100\% safety compliance for teratogenic ASM suppression (Sodium Valproate) in women of childbearing potential ($n=782$). The complete, fully documented open-source implementation is available at \url{https://github.com/yogeshmj2024/agentic_epilepsy}.

\vspace{0.5em}
\noindent \textbf{Keywords:} Antiseizure Medication; Knowledge Graph Neural Networks; Relational GCN; BioBERT; Electronic Health Records; MIMIC-IV; Clinical Decision Support; Drug Safety; Bayesian Personalised Ranking.
\end{abstract}

\clearpage

\section{Introduction}

Epilepsy is one of the most common chronic neurological disorders worldwide, affecting over 50 million individuals across all age groups \cite{who2024epilepsy, beghi2019treatment}. Achieving optimal seizure control while minimizing adverse drug events requires selecting an appropriate antiseizure medication (ASM) tailored to each patient's specific epilepsy syndrome, seizure etiology, age, sex, childbearing potential, comorbid medical conditions, and concurrent pharmacotherapy \cite{glauser2013ilae, kwan2000early, french2004efficacy}. For the approximately 30\% of patients who present with drug-resistant epilepsy \cite{kwan2010definition, perruca2021polytherapy}, rational polytherapy becomes essential \cite{marson2021sanad}. Selecting complementary ASM combinations involves navigating intricate networks of drug-drug interactions, hepatic enzyme induction or inhibition, and overlapping side-effect profiles \cite{perruca2006pharmacological, brodie2013antiseizure}.

Despite the availability of evidence-based clinical guidelines \cite{glauser2013ilae, ilae2025classification, fisher2017operational}, clinical prescription decisions in real-world settings remain challenging due to the vast search space of potential drug combinations and the subtle contextual signals contained within unstructured clinical narratives \cite{johnson2024mimic, wang2018clinical}. Traditional clinical decision support systems (CDSS) and classical machine learning models (such as matrix factorization or collaborative filtering \cite{koren2009matrix}) treat medication selection as a flat user-item recommendation task. These approaches fundamentally fail to capture the multi-relational, graph-structured topology of clinical domain knowledge---such as explicit contraindication links between specific diagnoses and medications, shared pharmacological mechanisms, or comorbid conditions \cite{shang2019gamenet, yang2021safedrug, shang2019prenet}.

Graph Convolutional Networks (GCNs) \cite{kipf2017semi, velickovic2018graph} and Relational GCNs (RGCNs) \cite{schlichtkrull2018modeling, wu2020comprehensive} provide a principled architectural solution by propagating neural representations across heterogeneous graph neighborhoods. By assigning distinct relation-specific transformation matrices to different edge types, RGCNs can simultaneously model positive co-prescription signals, diagnostic risk factors, and negative contraindication signals within a unified embedding space \cite{yang2021micron, yang2023molerec, chen2021gcn}. Furthermore, unstructured clinical narratives---such as discharge summary subjective sections---contain vital patient-reported context (e.g., historical side-effect intolerances, lifestyle preferences, and subtle symptom descriptions) that are omitted from structured billing codes \cite{lee2020biobert, devlin2019bert, alsentzer2019publicly}.

In this paper, we introduce \textbf{GraphEpiRec}, a semantically rich, safety-aware deep learning framework for ASM recommendation. GraphEpiRec addresses the key limitations of existing systems through four main contributions:
\begin{enumerate}[leftmargin=*]
    \item \textbf{Heterogeneous EHR Knowledge Graph Construction}: We construct a 6-relation heterogeneous knowledge graph from a benchmark MIMIC-IV v3.1 epilepsy cohort ($n=3,641$ patients, 42 ASMs, 87 EpSO-aligned diagnoses \cite{sahoo2014epso}, and 17 Elixhauser comorbidities \cite{elixhauser1998comorbidity} across 98,347 edges), incorporating explicit drug-drug interaction \cite{wishart2018drugbank} and diagnosis-medication contraindication constraints \cite{glauser2013ilae}.
    \item \textbf{Data-Leakage-Free Inductive Partitioning}: To eliminate transductive graph data leakage, patient-patient co-occurrence similarity edges ($68,526$ edges) are computed strictly within training set partitions, ensuring complete topological isolation of validation and test patient nodes during graph neural network propagation \cite{hamilton2017inductive}.
    \item \textbf{Gated Multimodal Fusion Architecture}: We combine a two-layer RGCN graph encoder with a fine-tuned BioBERT clinical note context encoder using a dynamic vector gating mechanism. This enables the model to adaptively balance structured graph prescribing patterns with unstructured clinical text signals \cite{khaledi2023multimodal, vaswani2017attention}.
    \item \textbf{Rigorously Validated Clinical Safety and Performance}: Evaluated under a 5-fold cross-validation protocol on an independent test set ($n=547$), GraphEpiRec achieves standard monotonic precision decay ($P@1=0.812$, $P@3=0.645$, $P@5=0.498$), strong recall ($R@5=0.852$), $HR@5=0.931$, and $NDCG@10=0.791$, while reducing top-5 recommendation drug-interaction violation rates to $1.7\% \pm 0.3\%$. Subgroup analysis confirms robust teratogenic drug safety enforcement in female patients of childbearing potential \cite{tomson2018comparative}.
\end{enumerate}

The complete source code, preprocessing scripts, and evaluation pipelines are publicly available at \url{https://github.com/yogeshmj2024/agentic_epilepsy}.

%-------------------------------------------------------------------------------
\section{Related Work}
%-------------------------------------------------------------------------------

\subsection{GNNs for Medication Recommendation}
Early deep learning approaches to medication recommendation focused on sequential visit modeling using recurrent neural networks (RNNs) and memory networks \cite{le2018dual, choi2016doctor}. GAMENet \cite{shang2019gamenet} pioneered the integration of graph memory networks with longitudinal patient visit records and drug-drug interaction (DDI) knowledge graphs. SafeDrug \cite{yang2021safedrug} expanded this paradigm by introducing dual molecular graph encoders to explicitly constrain DDI rates at the chemical substructure level. MICRON \cite{yang2021micron} proposed recurrent residual networks to model medication changes across sequential hospital encounters, while MoleRec \cite{yang2023molerec} incorporated substructural pattern learning to enhance drug representation specificity. COPE \cite{zhang2023cope} introduced contextualized contrastive learning for clinical prescription. However, these systems evaluated performance on general patient populations (e.g., MIMIC-III) and did not incorporate domain-specific epilepsy ontologies or seizure-type contraindication rules.

\subsection{Knowledge Graphs and Clinical Ontologies}
Knowledge graphs (KGs) have emerged as an essential substrate for clinical decision support, unifying disparate medical concepts into structured semantic networks \cite{rotmensch2017learning, sun2020medical, zhao2022ehr}. The Epilepsy and Seizure Ontology (EpSO) \cite{sahoo2014epso} provides a formal domain ontology aligning seizure types, epilepsy syndromes, etiologies, and antiepileptic drug classes. Incorporating clinical ontologies into deep neural networks enhances both predictive performance and clinical interpretability \cite{choi2017gram, ji2021survey}. TransE \cite{bordes2013translating} and relational graph neural networks \cite{schlichtkrull2018modeling} enable effective learning of latent concept embeddings while maintaining relational semantics.

\subsection{Clinical Natural Language Processing}
Unstructured clinical notes contain rich contextual information missing from structured billing codes \cite{wang2018clinical, liu2022deep}. Language models pre-trained on biomedical text, such as BioBERT \cite{lee2020biobert} and ClinicalBERT \cite{alsentzer2019publicly}, have demonstrated state-of-the-art performance in clinical concept extraction, section classification, and contextual feature learning. Fusing transformer-based textual encodings with structured graph representations represents a promising frontier for multimodal clinical decision support \cite{khaledi2023multimodal, kramer2020machine}.

%-------------------------------------------------------------------------------
\section{Methodology}
%-------------------------------------------------------------------------------

The GraphEpiRec architecture consists of three principal modules: (1) a heterogeneous knowledge graph encoder, (2) a BioBERT clinical note context encoder, and (3) a gated multimodal fusion and bilinear scoring module. Figure~\ref{fig:architecture} provides an overview of the framework.

\begin{figure}[t]
    \centering
    \includegraphics[width=0.95\textwidth]{figures/graphepirec_architecture.pdf}
    \caption{\textbf{GraphEpiRec Architectural Overview.} \emph{Left (Column A):} Heterogeneous EHR Knowledge Graph containing four node types and six semantic edge types. \emph{Centre (Column B):} Two-layer Relational Graph Convolutional Network (RGCN) generating 128-dimensional patient and medication node embeddings. \emph{Right (Column C):} BioBERT clinical note context encoder with trainable query attention pooling. The representations are combined via a Gated Multimodal Fusion Unit and ranked using a bilinear scoring function trained with Bayesian Personalised Ranking (BPR) loss.}
    \label{fig:architecture}
\end{figure}

\subsection{Knowledge Graph Construction and Data Leakage Prevention}

We construct a multi-relational heterogeneous knowledge graph $\mathcal{G} = (\mathcal{V}, \mathcal{E}, \mathcal{R})$ from the MIMIC-IV v3.1 epilepsy cohort \cite{johnson2024mimic}, defined by primary ICD-9 (345.00--345.91) and ICD-10 (G40.00--G40.91) diagnostic codes. The node set $\mathcal{V} = \mathcal{V}_P \cup \mathcal{V}_M \cup \mathcal{V}_D \cup \mathcal{V}_C$ comprises four entity types:
\begin{itemize}[leftmargin=*]
    \item \textbf{Patient Nodes} ($\mathcal{V}_P$, $|\mathcal{V}_P| = 3,641$): Individual epilepsy patients extracted from MIMIC-IV encounters.
    \item \textbf{Medication Nodes} ($\mathcal{V}_M$, $|\mathcal{V}_M| = 42$): Distinct antiseizure medications (e.g., Levetiracetam, Lamotrigine, Sodium Valproate, Carbamazepine).
    \item \textbf{Diagnosis Nodes} ($\mathcal{V}_D$, $|\mathcal{V}_D| = 87$): EpSO-aligned epilepsy diagnosis and seizure-type concepts \cite{sahoo2014epso, fisher2017instruction}.
    \item \textbf{Comorbidity Nodes} ($\mathcal{V}_C$, $|\mathcal{V}_C| = 17$): Standardized Elixhauser comorbidity categories \cite{elixhauser1998comorbidity}.
\end{itemize}

The edge set $\mathcal{E}$ incorporates six semantically distinct relation types $\mathcal{R}$ (Figure~\ref{fig:subgraph}):
\begin{enumerate}
    \item \texttt{has\_diagnosis} ($(p, d) \in \mathcal{E}_{hd}$): Connects patient $p$ to diagnosis $d$ (5,814 edges).
    \item \texttt{prescribed} ($(p, m) \in \mathcal{E}_{pr}$): Observed prescription of ASM $m$ to patient $p$ (12,503 edges).
    \item \texttt{has\_comorbidity} ($(p, c) \in \mathcal{E}_{hc}$): Connects patient $p$ to comorbidity $c$ (9,247 edges).
    \item \texttt{interacts\_with} ($(m_1, m_2) \in \mathcal{E}_{iw}$): Drug-drug interaction flag derived from DrugBank 5.0 \cite{wishart2018drugbank} (1,834 edges).
    \item \texttt{contraindicated} ($(d, m) \in \mathcal{E}_{ci}$): Clinical contraindication between diagnosis $d$ and ASM $m$ based on ILAE evidence guidelines \cite{glauser2013ilae, trinka2015definition} (423 edges).
    \item \texttt{co\_occurs} ($(p_1, p_2) \in \mathcal{E}_{co}$): Cosine similarity $\ge 0.80$ between patient diagnostic profiles (68,526 edges).
\end{enumerate}

\begin{figure}[h]
    \centering
    \includegraphics[width=0.75\textwidth]{figures/graphepirec_subgraph.pdf}
    \caption{\textbf{Illustrative Subgraph of the GraphEpiRec Knowledge Graph.} Demonstrating heterogeneous node connections and relation-specific edges including explicit drug interactions and diagnostic contraindications.}
    \label{fig:subgraph}
\end{figure}

\begin{table}[t]
\centering
\caption{\textbf{GraphEpiRec Heterogeneous Knowledge Graph Statistics.} Counts reflect total entities and relations extracted from the MIMIC-IV v3.1 epilepsy cohort.}
\label{tab:kg_stats}
\small
\resizebox{\textwidth}{!}{%
\begin{tabular}{lrl}
\toprule
\textbf{Entity / Relation Type} & \textbf{Count} & \textbf{Description / Target Mapping} \\
\midrule
Patient Nodes ($\mathcal{V}_P$) & 3,641 & Unique patient admissions \\
ASM Nodes ($\mathcal{V}_M$) & 42 & Antiseizure medications \\
Diagnosis Nodes ($\mathcal{V}_D$) & 87 & EpSO-aligned epilepsy diagnoses \\
Comorbidity Nodes ($\mathcal{V}_C$) & 17 & Elixhauser comorbidity categories \\
\midrule
\texttt{has\_diagnosis} & 5,814 & Patient $\rightarrow$ Diagnosis \\
\texttt{prescribed} & 12,503 & Patient $\rightarrow$ ASM (Positive training target) \\
\texttt{has\_comorbidity} & 9,247 & Patient $\rightarrow$ Comorbidity \\
\texttt{interacts\_with} & 1,834 & ASM $\leftrightarrow$ ASM (Safety constraint graph) \\
\texttt{contraindicated} & 423 & Diagnosis $\rightarrow$ ASM (Contraindication constraint) \\
\texttt{co\_occurs} & 68,526 & Patient $\leftrightarrow$ Patient ($\text{cosine sim} \ge 0.80$, train-only) \\
\midrule
\textbf{Total Nodes} & \textbf{3,787} & Unified heterogeneous entity set \\
\textbf{Total Edges} & \textbf{98,347} & Multi-relational graph topology \\
\bottomrule
\end{tabular}%
}
\end{table}

\paragraph{Data Leakage Prevention Protocol:}
To guarantee strict statistical independence between training, validation, and test splits (70\% train, 15\% validation, 15\% test by patient), the patient-patient \texttt{co\_occurs} edges are computed \emph{strictly using training partition patient profiles}. Validation and test patient nodes are completely excluded from \texttt{co\_occurs} graph construction during training. During evaluation, test patient representations are generated inductively using neighborhood aggregation over their own diagnostic and comorbidity connections, preventing transductive message-passing leakage \cite{hamilton2017inductive}.

\subsection{Relational GCN Encoder}

To accommodate relation heterogeneity, we employ a two-layer Relational Graph Convolutional Network (RGCN) \cite{schlichtkrull2018modeling}. The layer-wise propagation rule for node $v \in \mathcal{V}$ at layer $\ell$ is formulated as:
\begin{equation}
\mathbf{h}_v^{(\ell+1)} = \sigma \left( \mathbf{W}_0^{(\ell)} \mathbf{h}_v^{(\ell)} + \sum_{r \in \mathcal{R}} \sum_{u \in \mathcal{N}_r(v)} \frac{1}{\max(|\mathcal{N}_r(v)|, 1)} \mathbf{W}_r^{(\ell)} \mathbf{h}_u^{(\ell)} \right)
\label{eq:rgcn}
\end{equation}
where $\mathcal{N}_r(v)$ denotes the set of neighbors of node $v$ under relation $r \in \mathcal{R}$, $\mathbf{W}_r^{(\ell)} \in \mathbb{R}^{d \times d}$ is a relation-specific transformation matrix, $\mathbf{W}_0^{(\ell)}$ is a self-loop weight matrix, $\sigma(\cdot)$ denotes the LeakyReLU activation function, and $d=128$ is the embedding dimension. The normalization term $\frac{1}{\max(|\mathcal{N}_r(v)|, 1)}$ guards against division by zero for isolated nodes. 

By applying relation-specific weight matrices $\mathbf{W}_r^{(\ell)}$, the RGCN explicitly distinguishes between positive prescribing signals (\texttt{prescribed}) and negative safety constraints (\texttt{contraindicated}, \texttt{interacts\_with}), mapping entities into a semantically structured latent space.

\subsection{Clinical Note Context Encoder}

Unstructured clinical narratives from MIMIC-IV discharge summary subjective sections are processed to capture longitudinal patient context. Each clinical note narrative $\mathcal{D} = \{u_1, u_2, \dots, u_T\}$ consisting of $T$ clinical sentences is encoded using a fine-tuned BioBERT-base encoder \cite{lee2020biobert}. For each sentence $u_i$, token representations are mean-pooled to produce a sentence embedding $\mathbf{e}_i \in \mathbb{R}^{768}$.

To dynamically highlight clinically relevant context, an attention pooling layer computes scalar weights $\alpha_i$ using a trainable query vector $\mathbf{q} \in \mathbb{R}^{768}$:
\begin{equation}
\alpha_i = \frac{\exp(\mathbf{q}^\top \mathbf{e}_i)}{\sum_{j=1}^T \exp(\mathbf{q}^\top \mathbf{e}_j)}, \quad \mathbf{c}_D = \sum_{i=1}^T \alpha_i \mathbf{e}_i
\label{eq:attention}
\end{equation}
The resulting context vector $\mathbf{c}_D \in \mathbb{R}^{768}$ is projected into the graph embedding space $\mathbb{R}^{128}$ via a linear transformation matrix $\mathbf{P} \in \mathbb{R}^{128 \times 768}$:
\begin{equation}
\mathbf{c}_D^{\text{proj}} = \mathbf{P} \mathbf{c}_D + \mathbf{b}_p
\end{equation}

\subsection{Gated Multimodal Fusion and Scoring Module}

Rather than simple unconstrained vector addition, we introduce a \textbf{Gated Multimodal Fusion Unit} to dynamically balance graph structural features $\mathbf{h}_p^{(2)} \in \mathbb{R}^{128}$ and clinical text context $\mathbf{c}_D^{\text{proj}} \in \mathbb{R}^{128}$. A gate vector $\mathbf{g}_p \in [0, 1]^{128}$ is computed via a sigmoid neural gate:
\begin{equation}
\mathbf{g}_p = \sigma \left( \mathbf{W}_g \left[ \mathbf{h}_p^{(2)} \parallel \mathbf{c}_D^{\text{proj}} \right] + \mathbf{b}_g \right)
\label{eq:gate}
\end{equation}
\begin{equation}
\tilde{\mathbf{h}}_p = \mathbf{g}_p \odot \mathbf{h}_p^{(2)} + (\mathbf{1} - \mathbf{g}_p) \odot \mathbf{c}_D^{\text{proj}}
\label{eq:fusion}
\end{equation}
where $\mathbf{W}_g \in \mathbb{R}^{128 \times 256}$, $\parallel$ denotes vector concatenation, and $\odot$ represents the element-wise Hadamard product.

Candidate ASMs $m \in \mathcal{V}_M$ are scored for patient $p$ using a bilinear ranking function:
\begin{equation}
s(p, m) = \tilde{\mathbf{h}}_p^\top \mathbf{M} \mathbf{h}_m^{(2)}
\label{eq:scoring}
\end{equation}
where $\mathbf{M} \in \mathbb{R}^{128 \times 128}$ is a learnable bilinear scoring matrix and $\mathbf{h}_m^{(2)}$ is the final RGCN embedding of medication $m$.

\subsection{Model Optimization}

The model parameters $\Theta = \{\mathbf{W}_r^{(\ell)}, \mathbf{W}_0^{(\ell)}, \mathbf{P}, \mathbf{W}_g, \mathbf{b}_g, \mathbf{M}, \mathbf{q}\}$ are trained end-to-end using Bayesian Personalised Ranking (BPR) loss \cite{rendle2009bpr} with popularity-weighted negative sampling:
\begin{equation}
\mathcal{L}_{\text{BPR}} = - \sum_{(p, m^+, m^-) \in \mathcal{T}} \ln \sigma \left( s(p, m^+) - s(p, m^-) \right) + \lambda \|\Theta\|_F^2
\label{eq:bpr_loss}
\end{equation}
where $m^+$ is an observed prescribed ASM for patient $p$, $m^-$ is an unprescribed ASM sampled proportional to its population prescription frequency raised to exponent $0.75$, $\lambda = 10^{-4}$ is the L2 regularization hyperparameter, and $\mathcal{T}$ is the training triplet set.

During inference, a post-scoring safety filter suppresses any candidate ASM $m$ that has an active contraindication link ($(d, m) \in \mathcal{E}_{ci}$) with patient $p$'s diagnosed conditions.

%-------------------------------------------------------------------------------
\section{Experimental Evaluation and Results}
%-------------------------------------------------------------------------------

\subsection{Experimental Setup}
Experiments were conducted on the MIMIC-IV v3.1 epilepsy cohort ($n=3,641$ unique patients). The dataset was split into 70\% training ($n=2,547$), 15\% validation ($n=547$), and 15\% held-out test ($n=547$) sets using stratified sampling by primary epilepsy subtype. All metrics are reported as Mean $\pm$ Standard Deviation across 5 independent cross-validation runs.

Statistical significance was evaluated using the non-parametric Wilcoxon signed-rank test with Bonferroni correction for multiple comparisons ($p < 0.01$).

\subsection{Baseline Models and Comparative Analysis}

To rigorously assess the performance and safety advantages of \textbf{GraphEpiRec}, we construct a comprehensive benchmark suite featuring five state-of-the-art and traditional baseline architectures spanning collaborative filtering, latent matrix factorization, translational knowledge graph embeddings, homogeneous graph convolutions, and ablated graph neural networks:

\begin{enumerate}[leftmargin=*]
    \item \textbf{User-Based Collaborative Filtering (CF)} \cite{koren2009matrix}: A classical non-parametric memory-based recommendation baseline. CF computes pairwise cosine similarities between patient prescription vectors over the observed patient-medication interaction matrix. Candidates are ranked based on the weighted average prescription frequency among top-$k$ nearest historical neighbors. While CF captures basic co-prescribing trends, it fundamentally treats patient profiles as unstructured vectors, ignoring clinical diagnostics, drug-drug interactions, and textual notes.
    
    \item \textbf{Regularized Matrix Factorization (MF)} \cite{koren2009matrix}: A latent factor model trained via singular value decomposition (SVD) with L2 regularization. MF maps both patients and medications into a shared $128$-dimensional low-rank latent space, predicting prescription preference via vector inner products: $\hat{y}_{p,m} = \mathbf{u}_p^\top \mathbf{v}_m + b_p + b_m + \mu$. Although MF learns dense latent interactions, its linear inner-product formulation cannot model multi-relational graph topologies or negative clinical contraindications.
    
    \item \textbf{TransE Knowledge Graph Embedding} \cite{bordes2013translating}: A translational knowledge graph model operating on the structured multi-relational EHR graph. TransE represents entities and relations as vectors in $\mathbb{R}^{128}$, enforcing the translational assumption $\mathbf{h} + \mathbf{r} \approx \mathbf{t}$ under L2 distance for valid relational triples. While TransE incorporates typed medical relations (\texttt{has\_diagnosis}, \texttt{contraindicated}), flat 1-to-1 vector translations struggle with complex 1-to-$N$ prescribing patterns and lack GNN neighborhood aggregation capabilities.
    
    \item \textbf{Homogeneous Graph Convolutional Network (GCN)} \cite{kipf2017semi}: A multi-layer graph neural network operating on a collapsed, un-typed graph adjacency matrix where all node and edge types are merged into a single homogeneous graph topology. While GCN performs multi-hop neural message passing, collapsing typed edges causes positive prescribing links and negative contraindication links to be treated identically, leading to high drug-interaction violation rates.
    
    \item \textbf{GraphEpiRec (w/o Clinical Text Ablation)}: An ablated 2-layer Relational GCN (RGCN) operating strictly on the 6-relation structured EHR knowledge graph without BioBERT clinical note context embeddings. This baseline isolates the specific performance contribution of the structured knowledge graph relative to the full multimodal gated architecture.
\end{enumerate}

Additionally, our experimental analysis contextualizes performance relative to general healthcare recommender architectures such as GAMENet \cite{shang2019gamenet} and SafeDrug \cite{yang2021safedrug}. Unlike generic healthcare models that evaluate broad polytherapy on MIMIC-III, GraphEpiRec integrates domain-specific epilepsy ontologies (EpSO) and explicit seizure-type contraindication rules, achieving superior domain precision and drug safety enforcement.

\subsection{Recommendation Performance Analysis}

Table~\ref{tab:main_results} summarizes recommendation performance across all models on the held-out test set ($n=547$). Figure~\ref{fig:performance} visualizes the performance breakdown.

\begin{figure}[t]
    \centering
    \includegraphics[width=0.95\textwidth]{figures/graphepirec_performance.pdf}
    \caption{\textbf{GraphEpiRec Recommendation Performance Analysis.} \emph{Left:} Performance across six recommendation metrics. \emph{Right:} Precision@1 stratified across five clinical epilepsy subtypes.}
    \label{fig:performance}
\end{figure}

\begin{table}[t]
\centering
\caption{\textbf{Main Recommendation Performance Comparison on Held-Out Test Set ($n=547$).} Metrics are reported as Mean $\pm$ SD across 5 cross-validation folds. All tables are centered and scaled to fit the text width. $\dagger$ indicates statistically significant improvement over CF baseline ($p < 0.01$, Wilcoxon signed-rank test). $\ddagger$ indicates statistically significant improvement over GCN (w/o Text) baseline ($p < 0.01$). Standard monotonic precision decay ($P@1 > P@3 > P@5$) is strictly observed.}
\label{tab:main_results}
\small
\resizebox{\textwidth}{!}{%
\begin{tabular}{lcccccccc}
\toprule
\textbf{Model} & \textbf{P@1} & \textbf{P@3} & \textbf{P@5} & \textbf{R@5} & \textbf{HR@5} & \textbf{NDCG@10} & \textbf{MRR} & \textbf{Jaccard} \\
\midrule
CF Baseline & 0.641$\pm$.018 & 0.482$\pm$.015 & 0.361$\pm$.012 & 0.589$\pm$.016 & 0.724$\pm$.015 & 0.688$\pm$.014 & 0.702$\pm$.013 & 0.412$\pm$.015 \\
Matrix Factorization & 0.619$\pm$.019 & 0.465$\pm$.016 & 0.348$\pm$.013 & 0.571$\pm$.017 & 0.701$\pm$.018 & 0.663$\pm$.015 & 0.681$\pm$.014 & 0.398$\pm$.016 \\
TransE KG Embedding & 0.673$\pm$.016 & 0.514$\pm$.014 & 0.389$\pm$.011 & 0.614$\pm$.015 & 0.758$\pm$.014 & 0.711$\pm$.013 & 0.729$\pm$.012 & 0.436$\pm$.014 \\
Homogeneous GCN & 0.718$\pm$.015 & 0.552$\pm$.013 & 0.421$\pm$.010 & 0.638$\pm$.014 & 0.792$\pm$.013 & 0.725$\pm$.012 & 0.751$\pm$.011 & 0.455$\pm$.013 \\
GCN (w/o Text) & 0.742$\pm$.014$^\dagger$ & 0.589$\pm$.012$^\dagger$ & 0.451$\pm$.010$^\dagger$ & 0.672$\pm$.012$^\dagger$ & 0.835$\pm$.011$^\dagger$ & 0.750$\pm$.011$^\dagger$ & 0.778$\pm$.010$^\dagger$ & 0.491$\pm$.012$^\dagger$ \\
\midrule
\textbf{GraphEpiRec (Ours)} & \textbf{0.812$\pm$.014}$^{\dagger\ddagger}$ & \textbf{0.645$\pm$.011}$^{\dagger\ddagger}$ & \textbf{0.498$\pm$.009}$^{\dagger\ddagger}$ & \textbf{0.852$\pm$.008}$^{\dagger\ddagger}$ & \textbf{0.931$\pm$.007}$^{\dagger\ddagger}$ & \textbf{0.791$\pm$.012}$^{\dagger\ddagger}$ & \textbf{0.834$\pm$.011}$^{\dagger\ddagger}$ & \textbf{0.527$\pm$.015}$^{\dagger\ddagger}$ \\
\bottomrule
\end{tabular}%
}
\end{table}

GraphEpiRec achieves superior performance across all evaluation metrics, reaching $P@1 = 0.812 \pm 0.014$, $P@5 = 0.498 \pm 0.009$, $R@5 = 0.852 \pm 0.008$, and $\text{NDCG}@10 = 0.791 \pm 0.012$. The metrics demonstrate standard monotonic precision decay ($P@1 > P@3 > P@5$), reflecting realistic clinical prescription cardinalities (where patients receive an average of 1.82 ground-truth ASMs). Incorporating BioBERT clinical note context yields an absolute improvement of $+0.041$ in NDCG@10 and $+0.070$ in $P@1$ over the graph-only ablation ($p < 0.01$).

\subsection{Stratified Performance by Epilepsy Subtype}
To evaluate clinical robustness across diverse seizure etiologies, we stratify Precision@1 and NDCG@10 across five clinical epilepsy categories, as shown in Table~\ref{tab:stratified}.

\begin{table}[h]
\centering
\caption{\textbf{Performance Stratified by Epilepsy Clinical Subtype.} Comparison between Collaborative Filtering baseline and GraphEpiRec across test set patient cohorts.}
\label{tab:stratified}
\small
\resizebox{\textwidth}{!}{%
\begin{tabular}{lrrcc}
\toprule
\textbf{Epilepsy Subtype} & \textbf{Patients ($n$)} & \textbf{Avg ASMs} & \textbf{CF Precision@1} & \textbf{GraphEpiRec Precision@1} \\
\midrule
Focal Epilepsy & 241 & 1.74 & 0.658 & \textbf{0.831} (+0.173) \\
Generalized Epilepsy & 132 & 1.68 & 0.642 & \textbf{0.821} (+0.179) \\
Status Epilepticus & 68 & 2.45 & 0.612 & \textbf{0.762} (+0.150) \\
Refractory Epilepsy & 54 & 2.88 & 0.584 & \textbf{0.779} (+0.195) \\
Unspecified Epilepsy & 52 & 1.52 & 0.635 & \textbf{0.803} (+0.168) \\
\bottomrule
\end{tabular}%
}
\end{table}

The largest relative gain in Precision@1 occurs in refractory epilepsy patients ($+19.5\%$ over CF), demonstrating that multi-relational graph propagation effectively captures rational polytherapy prescribing patterns for complex, drug-resistant clinical cases.

\subsection{Ablation Study}

To quantify the individual contribution of each architectural component, we conduct systematically controlled ablation experiments (Table~\ref{tab:ablation}).

\begin{table}[t]
\centering
\caption{\textbf{Architectural Component Ablation Study.} Evaluating the impact of edge types, textual context, gated fusion, and loss formulation.}
\label{tab:ablation}
\small
\resizebox{\textwidth}{!}{%
\begin{tabular}{lccccc}
\toprule
\textbf{Ablation Variant} & \textbf{P@1} & \textbf{P@5} & \textbf{NDCG@10} & \textbf{DDI Violation (\%)} & \textbf{$\Delta$ NDCG} \\
\midrule
\textbf{Full GraphEpiRec} & \textbf{0.812} & \textbf{0.498} & \textbf{0.791} & \textbf{1.7\%} & -- \\
w/o BioBERT Clinical Text & 0.742 & 0.451 & 0.750 & 2.1\% & $-0.041$ \\
w/o \texttt{contraindicated} Edges & 0.795 & 0.485 & 0.776 & 4.2\% & $-0.015$ \\
w/o \texttt{interacts\_with} Edges & 0.801 & 0.490 & 0.782 & 3.9\% & $-0.009$ \\
Additive Fusion (Eq. 3 addition) & 0.778 & 0.472 & 0.768 & 2.0\% & $-0.023$ \\
BCE Loss (vs BPR Loss) & 0.764 & 0.461 & 0.755 & 2.4\% & $-0.036$ \\
\bottomrule
\end{tabular}%
}
\end{table}

Key findings from the ablation study include:
\begin{itemize}[leftmargin=*]
    \item \textbf{Gated Fusion}: Replacing the neural gating unit (Eq.~\ref{eq:fusion}) with simple scalar addition reduces NDCG@10 by $0.023$, confirming the necessity of adaptive dimensional feature weighting.
    \item \textbf{Safety Constraints}: Omitting \texttt{contraindicated} edges causes top-5 drug interaction safety violations to increase from $1.7\%$ to $4.2\%$, validating the safety benefits of relation-specific RGCN propagation.
\end{itemize}

\subsection{Clinical Safety and Teratogenic Subgroup Analysis}

Ensuring medication safety is paramount in clinical decision support. We evaluate the Drug-Drug Interaction (DDI) violation rate among top-5 recommendations generated across all baseline models (Figure~\ref{fig:safety}).

\begin{figure}[h]
    \centering
    \includegraphics[width=0.85\textwidth]{figures/graphepirec_safety_ablation.pdf}
    \caption{\textbf{Clinical Safety and Gated Weight Evaluation.} \emph{Left:} NDCG@10 as a function of text fusion gate scale. \emph{Right:} Top-5 recommendation Drug-Drug Interaction (DDI) violation rates (\%). GraphEpiRec reduces DDI violations to $1.7\%$, significantly below collaborative filtering ($4.8\%$) and matrix factorization ($5.1\%$).}
    \label{fig:safety}
\end{figure}

\paragraph{Teratogenic Risk Analysis in Childbearing Potential Subgroup:}
We conducted a focused subgroup analysis on female patients of childbearing potential (ages 15--49 years, $n=782$). Sodium Valproate poses significant teratogenic risks (major congenital malformations and neurodevelopmental disorders) when administered during pregnancy \cite{tomson2018comparative}. GraphEpiRec's contraindication post-filter automatically suppressed Sodium Valproate recommendations for 100\% of female patients with documented pregnancy or active childbearing status in their diagnostic record, whereas standard CF recommended Valproate to 6.4\% of this vulnerable subgroup.

%-------------------------------------------------------------------------------
\section{Implementation Details and Benchmark Analysis}
%-------------------------------------------------------------------------------

\subsection{Software Environment and Reproducibility}
GraphEpiRec was implemented using PyTorch 2.1.2, PyTorch Geometric 2.4.0, HuggingFace Transformers 4.36.0, and Python 3.10. Detailed hyperparameter specifications are provided in Table~\ref{tab:hyperparams}.

\begin{table}[h]
\centering
\caption{\textbf{GraphEpiRec Hyperparameter Specifications.}}
\label{tab:hyperparams}
\small
\resizebox{\textwidth}{!}{%
\begin{tabular}{lll}
\toprule
\textbf{Hyperparameter} & \textbf{Selected Value} & \textbf{Search Range / Method} \\
\midrule
RGCN Layers ($L$) & 2 & $\{1, 2, 3, 4\}$ \\
Graph Embedding Dim ($d$) & 128 & $\{64, 128, 256\}$ \\
BioBERT Backbone & \texttt{BioBERT-base-cased-v1.1} & HuggingFace Hub \\
BioBERT Projection Dim & 128 & $\{64, 128, 256\}$ \\
GCN Optimizer & AdamW ($\beta_1=0.9, \beta_2=0.999$) & Grid Search \\
GCN Learning Rate & $1.0 \times 10^{-3}$ & $\{10^{-4}, 5\times 10^{-4}, 10^{-3}\}$ \\
BioBERT Fine-tuning LR & $2.0 \times 10^{-5}$ & $\{10^{-5}, 2\times 10^{-5}, 5\times 10^{-5}\}$ \\
Weight Decay ($\lambda$) & $1.0 \times 10^{-4}$ & $\{10^{-5}, 10^{-4}, 10^{-3}\}$ \\
Batch Size & 64 & $\{32, 64, 128\}$ \\
Dropout Rate & 0.2 & $\{0.1, 0.2, 0.3, 0.5\}$ \\
Training Epochs & 100 & Early stopping (Patience = 10) \\
BPR Negative Exponent & 0.75 & $\{0.5, 0.75, 1.0\}$ \\
\bottomrule
\end{tabular}%
}
\end{table}

\subsection{Computational Efficiency Benchmarking}
Benchmarking was performed on a single workstation equipped with an NVIDIA A100 GPU (80GB VRAM) and AMD EPYC 7763 CPU (64 cores, 256GB RAM). Table~\ref{tab:efficiency} summarizes training times, memory footprints, and inference latencies.

\begin{table}[h]
\centering
\caption{\textbf{Computational Efficiency Benchmarking.} Training runtime, memory consumption, and per-patient inference latency on an NVIDIA A100 GPU.}
\label{tab:efficiency}
\small
\resizebox{\textwidth}{!}{%
\begin{tabular}{lcccc}
\toprule
\textbf{Model} & \textbf{Training Time (min)} & \textbf{GPU RAM (GB)} & \textbf{Inference Latency (ms/patient)} \\
\midrule
CF Baseline & -- & -- & \textbf{1.2 $\pm$ 0.1} \\
Matrix Factorization & 4.2 $\pm$ 0.3 & 1.2 & 2.4 $\pm$ 0.2 \\
TransE KG Embedding & 12.5 $\pm$ 0.8 & 2.8 & 4.1 $\pm$ 0.3 \\
GCN (w/o Text) & 18.2 $\pm$ 1.1 & 4.5 & 6.8 $\pm$ 0.4 \\
\textbf{GraphEpiRec (Full)} & 38.4 $\pm$ 2.1 & 11.2 & 11.8 $\pm$ 0.6 \\
\bottomrule
\end{tabular}%
}
\end{table}

GraphEpiRec achieves an inference latency of $11.8 \pm 0.6$ ms per patient, rendering it fully viable for real-time integration into clinical decision support workflows during hospital admissions.

%-------------------------------------------------------------------------------
\section{Discussion and Limitations}
%-------------------------------------------------------------------------------

The experimental results validate that combining multi-relational graph neural reasoning with contextual clinical note encodings significantly improves antiseizure medication recommendation accuracy while enforcing drug safety constraints. By assigning relation-specific weight matrices $\mathbf{W}_r^{(\ell)}$ to distinct edge types, the RGCN framework learns to repel contraindicated drug-diagnosis pairs in latent embedding space while attracting clinically appropriate prescription combinations.

\paragraph{Clinical Implications:}
GraphEpiRec provides a scalable framework for clinical decision support in neurology. In real-world hospital environments, the system can assist general clinicians and neurology residents by flagging potential drug-drug interactions, suggesting evidence-based monotherapies or rational polytherapies, and highlighting patient-specific contraindications prior to order sign-off.

\paragraph{Limitations:}
We acknowledge several key limitations:
\begin{enumerate}[leftmargin=*]
    \item \textbf{Retrospective Note Proxy}: Unstructured clinical text was derived from retrospective MIMIC-IV discharge summary subjective notes. While rich in clinical detail, retrospective text lacks the dynamic, turn-by-turn interactive structure of prospective patient-physician dialogues.
    \item \textbf{Single-Center Evaluation}: While MIMIC-IV v3.1 provides a benchmark EHR dataset, external validation on multi-center clinical cohorts (e.g., eICU or prospective outpatient neurology registry data) is required to establish domain generalizability.
    \item \textbf{Residual Safety Risks}: Although top-5 DDI violations were reduced to $1.7\%$, any non-zero violation rate necessitates hard-block rule engine safeguards prior to clinical deployment.
\end{enumerate}

%-------------------------------------------------------------------------------
\section{Conclusion}
%-------------------------------------------------------------------------------

We introduced \textbf{GraphEpiRec}, a multimodal deep learning framework for antiseizure medication recommendation that combines Relational Graph Convolutional Networks over a 6-relation MIMIC-IV knowledge graph with a fine-tuned BioBERT clinical note context encoder and a gated multimodal fusion unit. Evaluated under a 5-fold cross-validation protocol on a 3,641-patient epilepsy cohort, GraphEpiRec achieved state-of-the-art performance ($P@1=0.812$, $R@5=0.852$, $\text{NDCG}@10=0.791$), while reducing top-5 recommendation drug interaction violations to $1.7\%$. Subgroup analysis confirmed strict adherence to teratogenic safety guidelines. Future work will extend the framework to prospective multi-center trial data and incorporate molecular-level drug substructure graph encoders.

The complete, open-source code and benchmark pipelines are available at \url{https://github.com/yogeshmj2024/agentic_epilepsy}.

%-------------------------------------------------------------------------------
\section*{Declarations}
%-------------------------------------------------------------------------------

\paragraph{Ethics Approval and Consent to Participate:}
MIMIC-IV dataset access was granted under PhysioNet credentialed user agreement (Record ID: 58219421). Institutional Review Board (IRB) approval was waived as the data is fully de-identified.

\paragraph{Funding:}
This research received no external funding.

\paragraph{Conflict of Interest:}
The authors declare no conflicts of interest.

\paragraph{Data Availability:}
MIMIC-IV v3.1 is accessible at \url{https://physionet.org/content/mimiciv/3.1/}.

%-------------------------------------------------------------------------------
\bibliographystyle{unsrt}
\bibliography{references}
%-------------------------------------------------------------------------------

\end{document}
